\documentclass[12pt,a4paper]{article}
\usepackage[utf8]{inputenc}
\usepackage{amsmath}
\usepackage{amsfonts}
\usepackage{amssymb}
\usepackage{geometry}
\usepackage{booktabs}

\geometry{margin=1in}

\begin{document}

\begin{center}
    \large \textbf{The Mechanical Vacuum: \\ Frame-Dragging, Satiation Velocity, and the Variable Gravitational Constant} \\
    \vspace{1em}
    \normalsize Robin Skavberg \\
    \normalsize February 2026
\end{center}

\begin{abstract}
In standard General Relativity (GR), $G$ is a coupling constant that determines how much spacetime curves for a given amount of energy-momentum. GR treats $G$ as a static scalar, whereas this model treats it as a dynamic equilibrium state of the field itself. GR views $G$ in Einstein’s field equations\textsuperscript{1} (typically written) as:$$R_{\mu\nu} - \frac{1}{2}Rg_{\mu\nu} + \Lambda g_{\mu\nu} = \frac{8\pi G}{c^4}T_{\mu\nu}$$In this view, $G$ is an ``external'' value. It is not a property of the field that reacts to the tensor $T_{\mu\nu}$; it is a fixed ``conversion factor.'' GR does not consider $G$ to be in equilibrium with the tensor strength because $G$ is assumed to be a fundamental property of the universe's ``hardware,'' not a software-like variable that adjusts to local stress. \par
2. The Equilibrium Perspective (This Model) suggests that the ``constant'' $G$ is actually the restorative response of the vacuum field. If we define the tensor strength as the ``load'' on the vacuum, then $G$ must be in equilibrium with that load. Standard GR: $G$ is a constant. If $T_{\mu\nu}$ (matter) increases, the curvature ($R_{\mu\nu}$) increases linearly because the ``stiffness'' ($G$) is fixed. In this model: $G$ is a variable resultant. As $T_{\mu\nu}$ increases, the Vacuum Impedance increases, effectively ``damping'' or ``stretching'' the field's ability to respond. \par
3. Historically, Einstein did consider the idea that the inertia of a body depends on the presence of other masses (Mach's Principle), which would imply a variable gravitational influence. However, to keep the mathematics of General Relativity generally covariant and consistent with the Equivalence Principle, he settled on $G$ as a universal constant. We expand upon GR by suggesting that $G$ is in equilibrium with the tensor strength, the stiffness of spacetime is not uniform. It is ``softer'' where matter is dense (high impedance) and ``stiffer'' where matter is sparse (low impedance, high field tension). \par
This paper builds upon the original paper that described the mathematics of deriving the gravitational constant G as a function of vacuum restorative stiffness and matter-energy density\textsuperscript{2}:  $G = c^2 / (8\pi\rho_m)$. \par
We describe that the observed missing mass in galactic rotation curves is a result of the vacuum recovering its ground-state tension in low-density environments. We provide a mathematical description for field-strain satiation and propose a testable prediction regarding gravitational lensing discrepancies.
\end{abstract}

\section{The Mechanical Basis of Gravitation}
In standard relativistic astrophysics, G is treated as a universal constant. However, applying a mechanical framework to the vacuum suggests that G represents the coupling efficiency between matter and the elastic fabric of space. We define $c^2$ as the restorative stiffness of the vacuum medium. Matter-energy density $\rho_m$ acts as the mechanical load, leading to the derivation:
\begin{equation}
G = \frac{c^2}{8\pi\rho_m}
\end{equation}

\section{Impedance, Strain, and Field-Strain Satiation}
We define the vacuum as an elastic medium subject to impedance and strain. 

\subsection{Description of Terms}
\begin{itemize}
    \item $\rho_m$: Matter-Energy Density (The mechanical load).
    \item $c^2$: Restorative Stiffness (The intrinsic modulus of the vacuum).
    \item $Z_v$: Vacuum Impedance (Opposition to force).
    \item $\epsilon$: Field Strain (Deformation where $\epsilon \propto 1/G$).
\end{itemize}

\subsection{Mathematical Correlation}
In wave mechanics, impedance $Z_v$ is proportional to $\rho_m$. As $\rho_m$ increases, $Z_v$ rises, leading to \textbf{Strain Satiation}. This occurs when the field is mechanically saturated and the coupling G approaches its minimum. 

\section{Prediction for Galactic Rotation Curves}
The Dark Matter effect is reinterpreted as the recovery of the gravitational constant in low-impedance environments. Local measurements of G in high-density solar systems are suppressed. As $\rho_m$ drops in the galactic outskirts, G increases toward its ground-state maximum, providing the necessary centripetal force to maintain stellar orbits.

\section{Mechanical Validation and Scaling}
The transition from local suppressions to universal maxima can be categorized into three pillars of evidence:

\textbf{1. The Local Evidence (Mechanical Satiation)} \\
In high-density environments like the Earth's surface, the vacuum is "loaded" by mass. This creates a state of mechanical satiation where $G$ is at its minimum observed value ($6.67 \times 10^{-11}$). The resulting velocity $v_{\gamma} \approx 0.00004$ m/s matches the order of magnitude of linear frame-dragging (Lense-Thirring effect) measured by Gravity Probe B\textsuperscript{3}, identifying this velocity as the "viscosity" of the local vacuum frame.


\textbf{2. The Universal Floor (Unruh Satiation)} \\
The theoretical noise floor of the vacuum, as described by the Unruh Effect at $1g$\textsuperscript{4}, suggests a minimum thermal jitter. Our result of $40$ $\mu$m/s identifies the propagation floor of the vacuum. This "clamping" effect prevents the gravitational coupling from diverging in the presence of planetary mass, maintaining the Newtonian limit.


\textbf{3. The Galactic Solution (Dark Matter)} \\
The variance between the local satiation velocity and $c$ represents the dynamic range of the vacuum. As $\rho_m \to 0$ in galactic halos, the vacuum recovers its restorative stiffness, allowing $G$ to return to its ground-state maximum. This increased coupling provides the "extra" centripetal pull required to flatten rotation curves without the need for additional dark mass.


\section{Discussions}
This mechanical model predicts that lensing in low-density voids is enhanced. This hypothesis aligns with theories proposed by Christodoulou and Kazanas (2019)\textsuperscript{5}. Furthermore, studies of Void Lensing Anomalies (Sanchez et al., 2017)\textsuperscript{6} have noted discrepancies in lensing signals from certain cosmic voids, suggesting a possible variance in vacuum impedance.

\newpage
\appendix
\section{Technical Appendix}

\subsection{A1: Derivation of Velocity-Gravity Equivalence}
The core of this framework lies in the link between the gravitational coupling constant ($G$) and the propagation velocity of light ($v_{\gamma}$). 

\textbf{Step 1: Foundational Mechanical Expression} \\
We begin with G as a measure of field tension relative to density:
\begin{equation}
G = \frac{c^2}{8\pi\rho_m} \implies \rho_m = \frac{c^2}{8\pi G}
\end{equation}

\textbf{Step 2: Defining Vacuum Refractive Index} \\
Applying the Lorentz-Lorenz principle to the vacuum medium, the refractive index $n_v$ is determined by the mechanical load $\rho_m$:
\begin{equation}
n_v = \sqrt{1 + \rho_m}
\end{equation}

\textbf{Step 3: Deriving the Phase Velocity} \\
Substituting the expression for $\rho_m$ into the velocity equation $v_{\gamma} = c/n_v$:
\begin{equation}
v_{\gamma} = \frac{c}{\sqrt{1 + \frac{c^2}{8\pi G}}} = c \sqrt{\frac{8\pi G}{8\pi G + c^2}}
\end{equation}

\textbf{Step 4: Numerical Evaluation} \\
To evaluate the expression, we use standard SI values ($c \approx 2.9979 \times 10^8$ m/s, $G \approx 6.6743 \times 10^{-11}$ m$^3$ kg$^{-1}$ s$^{-2}$):
$$8\pi G \approx 1.677 \times 10^{-9}$$ 
$$c^2 \approx 8.987 \times 10^{16}$$ 
Plugging these into the ratio: $$\frac{8\pi G}{8\pi G + c^2} \approx 1.866 \times 10^{-26}$$ 
Taking the square root and multiplying by $c$: $$v_{\gamma} \approx 0.00004 \text{ m/s}$$ 

\subsection{Ground-State Restoration Table}
\begin{center}
\begin{tabular}{llll}
\hline
Environment & Density & Effective G & Rotation Effect \\ \hline
Solar System & High & G-standard & Newtonian \\
Galactic Disk & Medium & G $>$ G-standard & Partial Flatness \\
Galactic Halo & Low & G $\to$ G-max & Flat Rotation \\
Cosmic Void & $\to$ 0 & G-max & Max Lensing \\ \hline
\end{tabular}
\end{center}

\section*{Author Information}

\noindent \textbf{Author Contributions:} Robin Skavberg conceived of the mechanical vacuum model from his experience with electomagnetic wave theory, developed the mathematical derivation of G as a function of vacuum stiffness, and authored the manuscript. \\

\noindent \textbf{Competing Interests:} The author declares no competing interests. \\

\noindent \textbf{AI Statement:} Large Language Model (LLM) technology was utilized as a technical collaborator for LaTeX typesetting, document structure optimization, mathematical support, and bibliographic verification. The original unique physical theories, mathematical derivations, and scientific conclusions were produced entirely by the human author.

\begin{thebibliography}{9}
\bibitem{1} Einstein, A. (1916). Annalen der Physik, 354(7), 769--822.
\bibitem{2} Skavberg, R. (2026). The Mechanical Resolution of the Gravitational Constant. Zenodo. DOI: 10.5281/zenodo.18393277
\bibitem{3} Everitt, C. W. F., et al. (2011). Gravity Probe B: Test of General Relativity. Physical Review Letters, 106(22), 221101.
\bibitem{4} Unruh, W. G. (1976). Notes on black-hole evaporation. Physical Review D, 14(4), 870.
\bibitem{5} Christodoulou, D. M., \& Kazanas, D. (2019). Galactic rotation curves from a mechanical perspective. MNRAS.
\bibitem{6} Sanchez, C., et al. (2017). Cosmic voids and gravitational lensing anomalies. JCAP, 2017(09), 015.
\end{thebibliography}

\end{document}